\begin{examplebox}
	\textbf{\textcolor{CExample}{例 1（基础）}}
	\textbf{\textcolor{CExample}{题目：}}
	已知椭圆 $\dfrac{x^2}{4}+y^2=1$，点 $P\left(1,\dfrac{\sqrt{3}}{2}\right)$。求过点 $P$ 的切线方程。
	\textbf{\textcolor{CExample}{解题步骤：}}
	\begin{description}[style=nextline, leftmargin=2.8em, labelsep=0.5em, itemsep=0.45em]
		\item[\textbf{第一步（判断位置）：}]
			验证点 $P$ 在椭圆上：$\dfrac{1^2}{4}+\left(\dfrac{\sqrt{3}}{2}\right)^2=\dfrac14+\dfrac34=1$。
			\par\textbf{理由：} 点在曲线上，极线即为切线。
		\item[\textbf{第二步（代入极线）：}]
			极线方程为 $\dfrac{x_0x}{a^2}+\dfrac{y_0y}{b^2}=1$，代入 $x_0=1,y_0=\dfrac{\sqrt{3}}{2},a^2=4,b^2=1$ 得 $\dfrac{1\cdot x}{4}+\dfrac{\sqrt{3}/2\cdot y}{1}=1$。
			\par\textbf{理由：} 极线公式直接应用。
		\item[\textbf{第三步（化简结果）：}]
			化简得 $\dfrac{x}{4}+\dfrac{\sqrt{3}}{2}y=1$，即 $x+2\sqrt{3}y=4$。
			\par\textbf{理由：} 乘以最小公倍数化为整式。
	\end{description}
	\textbf{\textcolor{CExample}{关键结论：}}
	\textbf{切线方程为 $x+2\sqrt{3}y=4$。}

	\medskip

\textbf{例2：先判断切点是否实际存在。}
双曲线 $x^2-y^2/3=1$，$P=(2,1)$，求切点弦并判断可行性。
位置量 $q=4-1/3=11/3>1$，不满足两实切点条件 $q<1,\ q\ne0$。
其极线为 $2x-y/3=1$，即 $6x-y=3$。
代入曲线得 $11x^2-12x+4=0$，判别式 $144-176=-32<0$。
因此\textbf{实切点弦不存在}；$6x-y=3$ 只能称作极线，不能称作实切点连线。
\medskip
\textbf{\textcolor{CExample}{例 3（综合应用）}}
	\textbf{\textcolor{CExample}{题目：}}
	已知抛物线 $y^2=4x$，点 $A(0,1)$，点 $B$ 在抛物线上，且 $A$ 在 $B$ 的极线上。求点 $B$ 的坐标和极线方程。
	\textbf{\textcolor{CExample}{解题步骤：}}
	\begin{description}[style=nextline, leftmargin=2.8em, labelsep=0.5em, itemsep=0.45em]
		\item[\textbf{第一步（设极线）：}]
			设 $B(x_0,y_0)$，则极线方程为 $y_0y=2(x+x_0)$（$p=2$）。
			\par\textbf{理由：} 抛物线 $y^2=4x$ 对应 $2p=4$，故 $p=2$。
		\item[\textbf{第二步（代入条件）：}]
			由 $A$ 在极线上，有 $y_0\cdot1=2(0+x_0)$，即 $y_0=2x_0$。
			\par\textbf{理由：} 点满足极线方程。
		\item[\textbf{第三步（联立求解）：}]
			又 $B$ 在抛物线上：$y_0^2=4x_0$。代入 $y_0=2x_0$ 得 $(2x_0)^2=4x_0$，即 $4x_0^2=4x_0$，解得 $x_0=0$ 或 $1$，对应 $y_0=0$ 或 $2$。
			\par\textbf{理由：} 简单二次方程求解。
	\end{description}
	\textbf{\textcolor{CExample}{关键结论：}}
	\textbf{点 $B$ 坐标为 $(0,0)$ 或 $(1,2)$；对应极线方程分别为 $x=0$ 或 $y=x+1$。}
\end{examplebox}
